\documentclass[11pt,a4paper]{article}
\usepackage[margin=1in]{geometry}
\usepackage{graphicx}
\usepackage{hyperref}
\usepackage{enumitem}
\usepackage{titlesec}
\usepackage{fancyhdr}

% Header and footer
\pagestyle{fancy}
\fancyhf{}
\rhead{Week 3 Report}
\lhead{Online Trackers for Small Object Detection}
\rfoot{Page \thepage}

\title{\textbf{Weekly Progress Report - Week 3} \\
	\large Improve performance of existing online trackers for small object detection [UAV Videos]}

\author{Group: The Convolutioners\\\{Aagam Sheth, Mahima Parekh, Khevan Thanki, Riddhi Bhargava\}}

\date{Reporting Period: March 9 - March 14, 2026 \\
	Submission Date: March 14, 2026}

\begin{document}
	
	\maketitle
	
	\section*{Outline of Performed Tasks:}
	\begin{itemize}[leftmargin=*]
		\item Studied Siamese network architectures for object re-identification (ReID)
		\item Implemented a Siamese network model to learn similarity between object appearances
		\item Investigated contrastive loss and similarity metrics for identity matching
	\end{itemize}
	
	\section*{Siamese Network for Re-Identification}
	
	This week focused on exploring Siamese Networks as an approach for learning appearance similarity between objects. 
	A Siamese network consists of two identical neural network branches that share weights and process two input images simultaneously. 
	The network learns to determine whether the two inputs correspond to the same object identity.
	
	The model outputs feature embeddings for both inputs, and the similarity between these embeddings is measured using a distance metric such as Euclidean distance.
	This allows the system to determine whether two detections belong to the same tracked object.
	
	Siamese networks are particularly suitable for ReID tasks because they learn similarity relationships rather than direct class labels, making them effective for scenarios where new object identities may appear during tracking.
	
	\section*{Training Data Preparation}
	
	To train the Siamese network, image pairs were generated from object bounding boxes in the VisDrone dataset. The training samples consist of:
	\begin{itemize}
		\item \textbf{Positive pairs}: two object crops belonging to the same identity across different frames
		\item \textbf{Negative pairs}: object crops belonging to different identities
	\end{itemize}
	
	This pairing strategy enables the model to learn discriminative appearance representations that can distinguish between different tracked objects.
	
	\section*{Similarity Learning}
	
	The network is trained using a contrastive loss function, which encourages embeddings of the same object to be close together while pushing embeddings of different objects farther apart in the feature space.
	
	During inference, the learned embeddings can be used to compute similarity scores between detected objects and existing tracks in the tracking pipeline.
	
	This learned similarity measure helps improve identity association when motion cues alone are insufficient.
	
	\section*{Tentative List of Tasks for Next Week:}
	
	\begin{enumerate}
		\item Integrate the Siamese ReID module with the MOT tracking pipeline
		\item Evaluate similarity matching performance on VisDrone sequences
		\item Analyze identity switches and tracking consistency
	\end{enumerate}
	
\end{document}